\SourceLeaf{312}%
und die beigebrachten Zahlen enthalten diese Entwicklung bis zu den Grössen der achten Ordnung (einschl.), daher \(t,l,b\) selbst dadurch bis zu den Grössen der neunten Ordnung einschliesslich, oder der eilften Ordnung ausschliesslich bestimmt werden.

Die Entwicklung von \(\log b\) kann auch auf eine andere Art, nemlich nach den Potenzen von \(\beta,\tau\) und \(r\) geschehen. Setzt man
\[
z=\operatorname{tang}y
\]
so wird
\[
y=z-\tfrac13z^3+\tfrac15z^5-\tfrac17z^7+\tfrac19z^9-\text{u.\,s.\,w.}
\]
und hieraus
\[
\log y=\log z-\tfrac13zz+\tfrac{13}{90}z^4-\tfrac{251}{2835}z^6+\tfrac{3551}{56700}z^8-\text{u.\,s.\,w.}
\]
Wendet man diese Reihe an auf die Gleichung
\[
\operatorname{tang}\tfrac12b=\frac{\operatorname{tang}\frac12t}{\operatorname{tang}B\mathord{.}\operatorname{tang}T}
\]
nachdem man derselben vermöge der Substitutionen
\[
\begin{aligned}
\operatorname{tang}B\mathord{.}\operatorname{tang}T&=\frac\tau\beta\\
\operatorname{tang}\tfrac12t&=\frac{m\tau}{2\BareRoot{(1-\frac14mm\tau\tau)}}
\end{aligned}
\]
folgende Gestalt gegeben hat
\[
\frac{\beta m}{2\BareRoot{(1-\frac14mm\tau\tau)}}=\operatorname{tang}\tfrac12b
\]
so ergibt sich
\[
\begin{aligned}
\log b={}&\log\beta+\log m+\tfrac18mm\tau\tau+\tfrac1{64}m^4\tau^4+\tfrac1{384}m^6\tau^6+\tfrac1{2048}m^8\tau^8+\text{u.\,s.\,w.}\\
&-\tfrac1{12}\beta\beta(mm+\tfrac14m^4\tau\tau+\tfrac1{16}m^6\tau^4+\tfrac1{64}m^8\tau^6+\text{u.\,s.\,w.})\\
&+\tfrac{13}{1440}\beta^4(m^4+\tfrac12m^6\tau\tau+\tfrac3{16}m^8\tau^4+\text{u.\,s.\,w.})\\
&-\tfrac{251}{181440}\beta^6(m^6+\tfrac34m^8\tau\tau+\text{u.\,s.\,w.})\\
&+\tfrac{3551}{14515200}\beta^8(m^8+\text{u.\,s.\,w.})\\
&+\text{u.\,s.\,w.}
\end{aligned}
\]
oder, indem man die Glieder gleicher Ordnung zusammenfasst,
